% Chapter 7
\chapter{Innovative Design of a Handheld Surgical Instrument Balancing Weight and Functionality} % Main chapter title

\label{Chapter7} % For referencing the chapter elsewhere, use 

\section{Necessity To Develop Handheld Motorized Medical Instrument}
As established in prior discussions, contemporary research in surgical robotics has predominantly prioritized maximizing tip-positioning accuracy through advanced actuation methods and the integration of supplementary sensors. While these approaches demonstrably enhance precision, they often overlook a critical practical constraint: the weight-to-functionality ratio. The evident correlation between enhanced capabilities—achieved through more powerful motors or additional sensing—and increased instrument mass creates a significant adoption barrier. This added weight often leads surgeons to favor simpler, nearly weightless traditional instruments over their motorized counterparts, despite the latter's technological advantages.

To resolve this fundamental trade-off, the present work introduces a holistic approach to ergonomic and functional design in robotic-assisted surgery. This chapter details the conception, design, and implementation of a novel structured handheld instrument that strategically balances mass, functionality, and clinical usability. The proposed system embodies a dedicated effort to distribute system components intelligently, minimize inertial load, and integrate multi-modal sensing without compromising the manipulator's agility or the surgeon's comfort. By fundamentally re-evaluating the instrument architecture from a human-centric perspective, this research aims to deliver a platform that surgeons are not only able to use, but genuinely prefer to use, thereby bridging the gap between technological potential and clinical practicality.
%During the research and development of the Lasso Gripper, the functional principles of cables or ropes were extensively investigated, particularly in terms of their purpose and form when configured as a loop. Conventionally, a loop is formed by tying or securing a string, allowing object capture primarily through adjustments in the loop size. This method relies heavily on the motor's ability to execute rapid and precise grasping motions. Furthermore, the repeatability and consistency of such a system are strongly dependent on the design and performance of the launch and retraction mechanism.

%However, an innovative approach involves dividing the loop into two separate strands. By cutting the loop, two independent strings are created, each capable of applying tensile forces—pulling rather than pushing—on a target object. When these two strings are attached to opposite sides of an object and actuated in differentiated directions, they can generate in-plane rotational movements. This significantly enhances the dexterity and functionality of the system, enabling more complex manipulations such as rolling, reorienting, or stabilizing objects without requiring additional mechanical components.

%This concept is particularly relevant in the context of cable-driven hand-held medical devices. Such devices often require high precision, flexibility, and minimal invasive manipulation, especially in surgical or diagnostic applications. For instance, in laparoscopic tools, the use of dual-string actuation can allow for finer tip articulation and rotational control, improving the surgeon’s ability to navigate and interact with tissue in confined spaces. The tension-based control also offers inherent compliance and force feedback, which are critical for safe interaction with delicate biological structures. By utilizing the properties of the string, such medical instruments can both maintain a high force-volume ratio at the tip of the instrument while keeping a low profile in its size.

\section{Core Design Advantages for Hand-Held Applications}

This cable-driven, differential actuation philosophy offers several fundamental advantages that are particularly suited to the development of advanced hand-held devices:
    
    -\textbf{Enhanced Dexterity and Manipulation}: The system enables multi-degree-of-freedom manipulation (e.g., rotation, twisting) using a minimal number of actuated elements. This is crucial for applications requiring fine motor control within confined spaces.

    -\textbf{Inherent Compliance and Safety}: The system is naturally compliant. Cables under tension provide a direct force-feedback path. If an unexpected obstacle is encountered, the cables will slacken rather than apply a rigid, potentially damaging force. This passive safety feature is a critical asset for interactions in unstructured or delicate environments.

    -\textbf{Proximal Actuation and Minimal End-Effector Mass}: A primary benefit for hand-held device design is the ability to locate the heavy actuators (motors) and associated drive electronics proximally, either in the device's handle or even in a remote unit. This dramatically reduces the mass and inertia of the distal end-effector, leading to improved ergonomics, reduced user fatigue, and the potential for higher maneuverability and smaller form factors at the interaction point.

\section{Defining Performance Metrics for Structural Design}

To translate this conceptual framework into a viable product, key performance metrics must be established. These metrics will directly guide the structural design choices detailed in the subsequent chapter:

    -\textbf{Precision and Resolution}: The minimum achievable incremental movement of the end-effector tip, targeting sub-millimeter and sub-degree accuracy for delicate tasks.

    -\textbf{Force Capability}: The maximum tensile force the cables must transmit to securely grasp and manipulate objects of a target weight range.

    -\textbf{Range of Motion}: The required angular deflection and rotational range of the end-effector to perform its intended functions effectively.

    -\textbf{Force Transmission Efficiency}: Minimizing friction losses throughout the cable drive path is essential for accurate force feedback and predictable control.

    -\textbf{Low Hysteresis and Backlash}: The structural design must ensure minimal play and elastic deformation in the system to guarantee precise, repeatable, and immediate response to actuator inputs.

\section{Pathway to Structural Implementation}
This chapter has outlined the conceptual leap from a simple loop to a differentially actuated cable-driven system, highlighting its inherent advantages for dexterous manipulation in hand-held applications. The core principles of compliance, proximal actuation, and differential control form the foundation upon which a physical device must be built.

The following chapter will address the critical task of transforming this concept into a robust mechanical reality. The structural design must meticulously integrate several key subsystems:
\begin{enumerate}
    \item \textbf{Actuation Unit Selection and Layout}: Choosing motors and sensors that meet the performance metrics and determining their optimal arrangement.

    \item \textbf{Cable Routing and Guidance System}: Designing low-friction, precise pathways for the cables to ensure efficient force transmission and longevity.

    \item \textbf{End-Effector Configuration and Cable Termination}: Engineering the interface where the cables attach to the manipulator, ensuring secure anchoring, precise motion, and mechanical advantage.

    %\item \textbf{Ergonomic Handle and Housing Design}: Packaging the system into a form factor that is balanced, comfortable, and intuitive for the user to hold and operate over prolonged periods.
\end{enumerate}
Through this structured approach, the subsequent design will validate the conceptual advantages discussed herein and realize a functional, high-performance cable-driven hand-held device.

\subsection{Relation to the Lasso Gripper and Mechanism-Focused Thesis}

The design choices described in this chapter were strongly influenced by the earlier development of the Lasso Gripper. Rather than treating the handheld instrument as an isolated engineering project, this work positions the instrument as a second application that tests and extends mechanism-level hypotheses derived from the Lasso Gripper: that tensile elements can be used to increase capture adaptivity, that differential tendon routing produces compact multi-DOF behaviour, and that proximal actuation enables low distal inertia. Emphasising these commonalities keeps the thesis focused on theoretical and design principles for tendon-driven mechanisms, while using concrete device implementations to validate those principles.